\chapter{The AI Winter}

\section*{A Different Kind of Chapter}

The first two chapters each ended at a technical wall. Exhaustive symbolic reasoning ran into combinatorial explosion. The perceptron ran into a representational ceiling it couldn't cross with any amount of additional training. This chapter looks, on the surface, like a different kind of chapter, one about funding, reputation, and institutional support rather than mathematics. It isn't, and that's the point worth stating before the history: both of the collapses this chapter covers trace back to a specific, nameable technical limit from the first two chapters, not to a field simply falling out of fashion. The question this chapter answers is whether AI's history is a story of failed vision or a story of predictable technical bottlenecks arriving on schedule and colliding with promises that had outrun them. The evidence points firmly at the second.

It's tempting to read AI's history as a straight line: perceptron fails, the field waits, backpropagation arrives, deep learning follows. The actual sequence is messier and more useful to know, because the pattern it reveals, ambitious promises outrunning what the current bottleneck-resolution actually allows, then the funding retreat that follows, then a quieter period in which the next resolution gets worked out largely unnoticed, is not unique to the 1970s and 80s. It has recurred, in compressed form, several more times since, and later chapters will point back to this one when it does.

\section*{The First Winter}

Two reports did most of the damage.

In 1966, a US government committee called ALPAC evaluated a decade of heavily funded machine translation research, most of it aimed at automatically translating Russian scientific and military text for Cold War intelligence purposes. The committee's actual findings were more measured than their reception; ALPAC still saw value in machine-assisted translation and recommended continued work in computational linguistics generally. What survived in the public and political memory was a blunter message: fully automatic translation had not delivered on ten years of promises, and continued funding at the current level was hard to justify. American machine translation research went dark for most of the next two decades.

Three years later, Minsky and Papert's \emph{Perceptrons} supplied a second, more theoretical blow, the one Chapter 2 walked through directly. It didn't kill research into learning machines by itself, but it gave funders and skeptics a precise, mathematically airtight reason to doubt that the perceptron's simple learning rule was a path to anything more general. A representational ceiling, proven rather than merely observed, is a hard thing to argue with.

Then, in 1973, the British government commissioned applied mathematician Sir James Lighthill, who had no background in AI research, to assess the field for the Science Research Council. His report concluded that AI had failed to produce the major practical impact its proponents had promised, and it singled out combinatorial explosion, the same wall Chapter 1 described in \texttt{non\_sequitur.mac}, as a fundamental barrier the field had not found a way around. The British government responded by cutting AI funding across nearly every university department working on it. Researchers left the field or left the country. The report's influence didn't stop at the English Channel; it gave American funding skeptics ammunition too, and in 1974 DARPA, which had been AI research's largest and most patient funder through the 1960s, began sharply reducing its support.

By the mid-1970s, AI research in both countries had lost most of its institutional funding, and "artificial intelligence" had become, for a working researcher, a label that made it harder rather than easier to get a grant approved. This period, roughly 1974 to 1980, is what's usually meant by the first AI winter.

\section*{An Irony Worth Naming Early}

Here is a detail that rarely makes it into the shorter tellings of this story, and it's worth sitting with before moving on, because Chapter 4 depends on it. The mathematical technique that would eventually let networks with more than one layer be trained efficiently, resolving exactly the representational ceiling Minsky and Papert had proven, was first worked out, in close to its modern form, by Paul Werbos in his 1974 Harvard doctoral thesis. That is the same year DARPA began its retrenchment. The solution to the problem that had helped freeze the field's funding existed in a dissertation on a library shelf while the freeze was setting in. It would take more than a decade, and a different set of researchers working largely independently, before the technique was rediscovered, popularized, and connected back to neural networks widely enough to matter. Funding cycles and technical progress do not always move together, and this is the first of several points in this book's history where the useful idea was already sitting somewhere, waiting for the field's attention to come back around to it.

\section*{The Second Boom and the Second Winter}

AI funding and enthusiasm did recover, but not by returning to neural networks. The 1980s boom was built on the expert systems described in Chapter 1. XCON, deployed at Digital Equipment Corporation starting in 1978 to configure VAX computer orders, was the clearest commercial success story the field had yet produced, reportedly saving the company tens of millions of dollars a year by encoding the expertise needed to translate a customer order into a correct, buildable system configuration. Stories like that one drove real investment. Companies built and sold expert system "shells," specialized Lisp machines optimized to run this kind of software commanded premium prices, and by the mid-1980s the AI industry had grown from a niche research pursuit into a market worth well over a billion dollars.

It did not last. Two separate problems converged around 1987. The first was economic: general-purpose workstations from Apple, Sun, and IBM caught up to and then undercut the price and performance of specialized Lisp hardware, and the market for that hardware collapsed within a couple of years, taking several companies down with it. The second was the same knowledge acquisition bottleneck from Chapter 1, now showing up at commercial scale. Expert systems built from thousands of rules were expensive to maintain, could not learn from new cases on their own, and broke, often badly, the moment a situation fell outside the narrow domain their rules had anticipated. Companies that had invested heavily discovered the maintenance cost of keeping a large rule base current, in constant collaboration with expensive human domain experts, exceeded the value the system provided. DARPA cut funding again. The second AI winter, running roughly from 1987 into the early 1990s, was in some ways deeper and more discouraging than the first, because it followed a period of genuine, well-documented commercial success rather than merely unmet promises.

\section*{What This Chapter Sets Up}

Two winters, separated by about a decade, driven by two different technologies and two different failure modes, funding collapse after unmet promises the first time, funding collapse after real but unmaintainable success the second time. It would be easy to read this as a field that simply couldn't get out of its own way. The more accurate reading is narrower and more specific: both winters followed directly from one of the two bottlenecks named in the first three chapters of this book, computational explosion in the first case, knowledge acquisition in the second, colliding with expectations that had outrun what either paradigm could actually deliver at scale.

Neither winter meant research stopped entirely. Small groups kept working on neural networks through the years when the label was most unfashionable, refining exactly the technique Werbos had described in 1974. Chapter 4 picks up there, with the paper that finally brought that technique to the field's attention and showed, concretely, that a network with more than one layer could learn to solve problems a single-layer perceptron provably never could.
